% Nasza inżynierka
\documentclass[10pt,twoside,a4paper,hidelinks]{article}
\usepackage{polski}                   % żeby były polskie napisy
\usepackage{fontspec}                 % żeby ustawić czcionkę na systemową (Arial)
\usepackage{geometry}                 % do marginesów
\usepackage{indentfirst}              % żeby pierwsze akapity też miały wcięcie
\usepackage{titlesec}                 % żeby formatować tytuły rozdziałów itd.
\usepackage{secdot}                   % aby dodać kropkę za numerkiem podrzodziałów i podpodrozdziałów
\usepackage{chngcntr}                 % umożliwia numerowanie obrazków itp. względem rozdziału
\usepackage{tocloft}                  % umożliwia ustawienia dotyczące spisu treści i innych spisów
\usepackage{tabu}                     % do tabel
\usepackage[table,dvipsnames]{xcolor} % do kolorowania tabel
\usepackage{tabularx}                 % do lepszych tabel
\usepackage[backend=biber,language=polish]{biblatex} % do bibliografii
\usepackage{enumitem}                 % do modyfikacji listy (begin{itemize}, niepotrzebne odstępy przed i po)
\usepackage{floatrow}                 % aby umożliwić wymuszenie położenia figury
\usepackage{caption}                  % do zmiany podpisów tabel i obrazków
\usepackage{setspace}                 % również do zmiany podpisów (konkretniej interlinii w podpisach - wymagany przez caption)
\usepackage{longtable}
\usepackage{booktabs}
\usepackage{tabu}
\usepackage{changepage}
\usepackage{amsthm}
\usepackage{hyperref}
\usepackage[final]{pdfpages}
\usepackage{multirow}
\usepackage[titletoc,title]{appendix}
\usepackage{listings}
\usepackage{dirtree}
\usepackage{multicol}
\usepackage{subcaption}
\usepackage{afterpage}
\usepackage{pgfplots}
\usepackage{pgfplotstable}
\usepackage{lipsum} % na potrzeby szablonu

\usepackage{fourier} 
\usepackage{array}
\usepackage{makecell}

\setmainfont{Arial}

\newcommand\myicon[1]{{\color{#1}\rule{2ex}{2ex}}}
\newcommand{\myfolder}[2]{\myicon{#1}\ {#2}}

\definecolor{listinggray}{gray}{0.9}
\definecolor{lbcolor}{rgb}{0.9,0.9,0.9}
\lstset{
  backgroundcolor=\color{lbcolor},
  tabsize=4,
  language=C++,
  captionpos=b,
  tabsize=4,
  frame=lines,
  numbers=left,
  numberstyle=\tiny,
  numbersep=5pt,
  breaklines=true,
  showstringspaces=false,
  basicstyle=\linespread{0.94}\ttfamily,
%  identifierstyle=\color{magenta},
  keywordstyle=\color[rgb]{0,0,1},
  commentstyle=\color{OliveGreen},
  stringstyle=\color{red}
  }


% Magia żeby dodatki miały ma początku "Dodatek X: ..."
\makeatletter
\renewcommand{\@redotocentry@pp}[1]{%
  \let\oldacl@pp=\addcontentsline
  \def\addcontentsline##1##2##3{%
    \def\@pptempa{##1}\def\@pptempb{toc}%
    \ifx\@pptempa\@pptempb
      \def\@pptempa{##2}\def\@pptempb{#1}%
      \ifx\@pptempa\@pptempb
        \oldacl@pp{##1}{##2}{\appendixname\space ##3}%
      \else
        \oldacl@pp{##1}{##2}{\chaptertitlename\space ##3}% added \chaptertitlename
      \fi
    \else
      \oldacl@pp{##1}{##2}{##3}%
    \fi}
}
\makeatother

% to jest jakas magia, aby odczytac szerokość longtable i ustawić tę wartość później jako LTcapwidth (parametr kontrolujący szerokość captiona w longtable)
% Creditsy i flaszkę proszę wysyłać do magika Heiko (http://compgroups.net/comp.text.tex/longtable-tablewidth/1922986)
\makeatletter
\newlength\LongtableWidth
\newcommand*{\org@longtable}{}
\let\org@longtable\longtable
\def\longtable{%
  \begingroup
    \advance\c@LT@tables\@ne
    \edef\x{LT@\romannumeral\c@LT@tables}%
    \global\LongtableWidth\z@
    \@ifundefined{\x}{%
      % longtable width not available
    }{%
      \def\LT@entry##1##2{%
        \global\advance\LongtableWidth##2\relax
      }%
      \@nameuse{\x}%
    }%
    % debug output
    \typeout{* \x: \the\LongtableWidth}%
  \endgroup
  \ifdim\LongtableWidth>\z@
    \setlength{\LTcapwidth}{\LongtableWidth}%
  \fi
  \org@longtable
}
\makeatother

\newcommand\blankpage{%
    \null
    \thispagestyle{empty}%
    \addtocounter{page}{-1}%
    \newpage}

% Rozpocznij od nowa numerowanie rysunków dla każdego rozdziału (section).
% Dodaje numer rozdziału do numeru rysunku: nr_rozdzial.nr_rysunku_w_ramach_rozdzialu
%
% Źródło: http://tex.stackexchange.com/questions/28333/continuous-v-per-chapter-section-numbering-of-figures-tables-and-other-docume
\counterwithin{figure}{section}

% to samo dla tabel
\counterwithin{table}{section}

% żeby nie było Rysunek tylko Rys
\renewcommand{\figurename}{Rys.}

% żeby nie było odstępów przed/po/w środku listy (itemize, ew. dodać też dla enumerate?)
\SetLabelAlign{parright}{\parbox[t]{\labelwidth}{\raggedleft#1}}
\setlist[itemize]{noitemsep,nolistsep,topsep=0pt}
\setlist[description]{noitemsep,nolistsep,topsep=0pt,style=multiline,align=parright}
    
% ustawiamy domyślny odstęp przed i po pływającymi elementami (tabele i obrazki) umieszczonymi w środku tekstu (flaga H) na 0
\setlength{\intextsep}{0mm}
\setlength{\textfloatsep}{0mm}

% ustawiamy domyślną czcionkę dla podpisów na small (9pt dla article 10pt) oraz interlinię na 1.0
\captionsetup{font={small,stretch=1.5}}

% własny separator do podpisów (to, co jest po 'Rys. X.Y' - kropka, nazwana wewnętrznie jako 'dot')
\DeclareCaptionLabelSeparator{dot}{. }
\DeclareFloatVCode{6ptskip}{\vspace{6pt}}
\DeclareFloatVCode{12ptskip}{\vspace{12pt}}

% dla tabel: 0pt pod opisem, 6pt nad
\captionsetup[table]{singlelinecheck=false} % nie wyśrodkowywuj opisu w pojedynczej linii
\captionsetup[table]{labelfont=bf,labelsep=dot} % pogrubienie nagłówka podpisu (Tabela X.Y) i zakonczenie jej wczesniej zdefiniowana kropką
\floatsetup[table]{font={small,stretch=1.0},capposition=top,captionskip=0pt,precode=12ptskip,postcode=12ptskip} % nie wiem dlaczego, aby otrzymac odstep 6pt przed tabela, trzeba tutaj dac 12pt :/

\setlength{\LTpre}{12pt}
\setlength{\LTpost}{12pt}

\tabulinesep=2.0mm % wzięte z czapy, ale wygląda dobrze (minimalny odstęp między początkkiem i końcem wierwsza a jego zawartością - przydaje się w przypadku zawijanych wierszy)

% dla obrazków: 6pt nad opisem, 12pt pod
\captionsetup[figure]{justification=centering} % inaczej niz w tabelach - zawsze centruj podpis
\captionsetup[figure]{labelsep=dot} % użyj kropki jako separatora ale nie pogrubiaj
\floatsetup[figure]{capposition=bottom,captionskip=6pt,precode=12ptskip,postcode=12ptskip}

% koment do poniższych: bfseries oznacza pogrubienie, itshape kursywę, mdseries normalną
% large = 12pt, small = 9pt (zależne od ustawionego u góry podstawowego 10pt),
% normalsize podstawowy rozmiar czyli 10pt

% formatowanie tytułów rozdziałów (tutaj nazwane sekcjami)

\makeatletter
\newcommand{\setappendix}{Dodatek~\thesection:}
\newcommand{\setchapter}{\thesection.}
\titleformat{\section}{\bfseries\large}{%
  \ifnum\pdfstrcmp{\@currenvir}{appendices}=0
    \setappendix
  \else
    \setchapter
  \fi}{0.5em}{\MakeUppercase}
\makeatother

	
% formatowanie tytułów podrozdziałów (tutaj nazwane podsekcjami)
\titleformat*{\subsection}{\normalsize\bfseries\itshape}
\sectiondot{subsection}

% formatowanie tytułów punktów podrozdziałów (tutaj nazwane podpodsekcjami)
\titleformat*{\subsubsection}{\normalsize\itshape}
\sectiondot{subsubsection}

\titlespacing*{\section}{0pt}{12pt}{6pt}
\titlespacing*{\subsection}{0pt}{4pt}{2pt}
\titlespacing*{\subsubsection}{0pt}{4pt}{2pt}

% SPIS TREŚCI

% żeby w spisie treści były też kropki po numerkach rozdziałów i podrozdziałów itd.
\renewcommand{\cftsecaftersnum}{.}
\renewcommand{\cftsubsecaftersnum}{.}
\renewcommand{\cftsubsubsecaftersnum}{.}
% żeby napis SPIS TREŚCI był wielkimi literami, pogrubiony itd
\renewcommand{\cfttoctitlefont}{\normalfont\large\bfseries\MakeUppercase}
% żeby tytuły rozdziałów w spisie oraz numery stron nie były pogrubione
\renewcommand\cftsecfont{\normalfont}
\renewcommand\cftsecpagefont{\normalfont}
% żeby dla rozdziałów też były kropki od napisu do numeru strony
\renewcommand\cftsecleader{\cftdotfill{\cftdotsep}}
% odstępy między akapitami 6pt
\setlength\cftbeforesecskip{6pt}
\setlength\cftbeforesubsecskip{6pt}
\setlength\cftbeforesubsubsecskip{6pt}
% żeby kropki od napisu do numeru strony były gęstsze
\renewcommand{\cftdotsep}{0}

% SPISY
% zmiana nazwy, czcionki, marginesu i separatora dla listy figur
\renewcommand{\listfigurename}{Wykaz rysunków}
\renewcommand{\cftloftitlefont}{\normalfont\large\bfseries\MakeUppercase}
\setlength\cftbeforefigskip{6pt}
\renewcommand{\cftfigaftersnum}{.}

\setlength\cftfignumwidth{1.1cm}

% zmiana nazwa, czcionki, marginesu i separatora dla listy tabel
\renewcommand{\listtablename}{Wykaz tabel}
\renewcommand{\cftlottitlefont}{\normalfont\large\bfseries\MakeUppercase}
\setlength\cftbeforetabskip{6pt}
\renewcommand{\cfttabaftersnum}{.}

\setlength\cfttabnumwidth{1.1cm}

% zmiana nazwy z 'Bibliografia' na 'Wykaz literatury'
\DefineBibliographyStrings{polish}{references = {Wykaz literatury}}
\DeclareFieldFormat{labelnumberwidth}{#1\adddot\midsentence}

% ustawienie marginesów
\geometry{
 a4paper,
 inner=3.5cm,
 outer=2.5cm,
 top=2.5cm,
 bottom=2.5cm
 }

\setlength{\parindent}{1.25cm}          % wcięcie przed akapitem
\renewcommand{\baselinestretch}{1.5}    % interlinia
\setlength{\parskip}{0pt}               % odległość pomiędzy akapitami

% żeby nie było wdów i sierot (linii samotnych ale nie słów!)
\widowpenalty10000
\clubpenalty10000

% komenda ignorująca cos w srodku tekstu (do naszych komentarzy)
\newcommand{\ignore}[2]{\hspace{0in}#2}

\theoremstyle{definition}
\newtheorem{definition}{Definicja}[section]

\theoremstyle{remark}
\newtheorem*{remark}{Komentarz}

\addbibresource{meta/bibliografia.bib}
\makeindex

\begin{document}   % właściwy początek dokumentu

\includepdf{meta/StronaTytulowa_155170.pdf}
\includepdf{meta/Oswiadczenie_155170.pdf}

\setcounter{page}{3} % rozpoczecie od strony 3
\section*{Streszczenie}
\lipsum[5]

\vspace{1cm}
\noindent
\textbf{Słowa kluczowe:} słowa kluczowe

\vspace{1cm}
\noindent
\textbf{Dziedziny nauki i techniki zgodne z wymogami OECD:} 1.2 Nauki o komputerze i informatyce
 \newpage
\section*{Abstract}
\lipsum[5]

\vspace{1cm}
\noindent
\textbf{Keywords:} keywords


\vspace{1cm}
\noindent
\textbf{OECD field of science and technology:} 1.2 Computer and information sciences
 \newpage
\clearpage
\renewcommand{\baselinestretch}{1.0}\normalsize	% żeby w spisie treści intelinia była pojedyncza
\tableofcontents
\renewcommand{\baselinestretch}{1.5}\normalsize \newpage
% http://tex.stackexchange.com/questions/30122/generate-table-of-contents-when-section-sections-without-numbering-has-been

\clearpage
\section*{Wykaz ważniejszych oznaczeń i skrótów} % section* - ukrywa numerowanie oraz wyklucza ze spisu tresci
\addcontentsline{toc}{section}{Wykaz ważniejszych oznaczeń i skrótów}% % reczne dodanie do spisu tresci

\begin{description}[leftmargin=2.5cm,labelwidth=2cm]
\item[PG] -- Politechnika Gdańska
\item[WETI] -- Wydział Elektroniki, Telekomunikacji i Informatyki
\item[CPU] -- z ang. Centrall Processing Unit - główny procesor urządzenia
\end{description}
 \newpage

%tutaj rozdzialy
\section{Wprowadzenie}

\subsection{Cel pracy}

 \newpage
\section[Wstęp do Api OpenCL]{Wstęp do Api OpenCL}
OpenCL czyli Open Compute Language to standard tworzony obecnie przez grupę Khronos, służący do pisania programów, które mogą zostać wykonane na różnych platformach takich jak CPU, GPU czy FPGA. Specyfikacja OpenCL definiuje interfejs w języku c++, który umożliwia zaprogramowanie aplikacji by ta wykonała konkretny kod na wybranym urządzeniu. Standard OpenCL jest głównie wykorzystywany do równoległych obliczeń takie jak wektorowe operacje matematyczne czy przetwarzanie obrazów. Za implementacje sterownika który wystawia api zgodne z określoną wersją specyfikacji odpowiedzialny jest producent urządzenia.  Dzięki temu, że standard jest otwarty a jego implementacje posiada większość producentów możemy stworzyć kod który możemy uruchomić niezależnie od architektury czy producenta posiadanego procesora głównego czy graficznego. Jest to duża zaleta w porównaniu na przykład do CUDA, która to jest interfejsem implementowanym jedynie przez Nvidie. Standard OpenCL jest rozwijany i modyfikowany, przez co api zdefiniowane jest w kilku wersjach. Najnowsza wersja specyfikacji to wersja 3.0. Wszystkie wersje są kompatybilne wstecznie. Dodatkowo zdefiniowane są rozszerzenia api takie jak cl\_khr\_gl\_sharing czy definiujący api do sharingu obiektów między OpenCL a OpenGL, takie dodatkowe api jest też specyfikowane przez grupę Khonosa w ramach określonej wersji OpenCL, jednak nie jest obowiązkowe. Istnieją także rozszerzenia api wyspecyfikowane przez konkretnego producenta np cl\_intel\_mem\_force\_host\_memory, które jest dostępna na urządzeniach intela, lub cl\_qcom\_android\_native\_buffer\_host\_ptr dostępny na procesorach Qualcoma z systemem Android. Takie dodatkowe api uzupełnia podstawę, umożliwiając lepsze dopasowanie specyfikacji do konkretnego sprzętu. OpenCL na platformie Android dostępny jest jedynie z poziomu natywnej biblioteki w języku c++.

\label{chap:goal}
\subsection{Linkowanie biblioteki OpenCL na platformie Android}
Aplikacje zwykle nie linkują się bezpośrednio ze sterownikiem posiadającym kompletną implementację api. Do tego wykorzystywana jest dodatkowa biblioteka która wyszukuje implementacji sterownika dla wszystkich platform na urządzeniu. Dzięki wykorzystaniu takiej ładującej biblioteki Aplikacja może używać każdej dostępnej platformy wspierającej to Api, oraz nie jest na sztywno połączona z jednym sterownikiem w określonej wersji.
Niestety nie istnieje binarna wersja takiej biblioteki dla systemu Android. Istnieje możliwość połączenia natywnej biblioteki ze sterownikiem znajdującym się w urządzeniu, jednak wadą takiego rozwiązania jest konieczność pobrania z urządzenia pliku binarnego z implementacją OpenCL oraz wszystkich zależnych od niej sterowników. Takie rozwiązanie powoduje, że skompilowana aplikacja będzie działać jedynie na urządzeniu, z którego zostały pobrane biblioteki.
Innym rozwiązaniem jest pobranie źródeł z kodem biblioteki ładującej, zbudowanie biblioteki i połączenie jej z aplikacją. Sterownik który będzie łączył program z implementacją OpenCL ma zapisane domyślne ścieżki w których mogą znajdować sie biblioteka OpenCL. Istnieje również możliwość zdefiniowania zmiennej środowiskowej w której wskazujemy lokalizację z której chcemy by sterownik został załadowany. Wadą takiego rozwiązania jest konieczność kompilacji takiej biblioteki, natomiast dzięki temu możemy zbudować aplikację działającą na wielu urządzeniach.
https://github.com/krrishnarraj/libopencl-stub

\subsection{Typowy przebieg aplikacji OpenCL}
Poniżej widać jak wygląda przebieg prostego programu wykoystującego Api OpenCL do inkrementowania każdego elementu bufora pamięci.

\lstinputlisting{listings/sample_cl.cpp}

Pierwszym krokiem jest zawołanie \textbf{clGetPlatformIDs}. Podając drugi argument czyli cl\_platform\_id\* jako null, do trzeciego argumentu jakim jest cl\_uint\* zostanie wpisana liczba dostępnych na urządzeniu platform wspierajacych OpenCL.
Drugie wywołanie clGetPlatformIDs wpisze informacje o podanej liczbie platform i zapisze je w tablicy podanej w drugim argumencie.
W tej pracy w testowanych aplikacjach na system android wykorzystywany będzie CL\_DEVICE\_TYPE\_GPU, który jako jedyny jest dostępny na testowanych przeze mnie urządzeniach.
Posiadając obiekt device, wołając \textbf{clCreateContext} możemy stworzyć context w ramach którego możliwe jest zarządzanie później stworzonymi obiektami na określonych przy tworzeniu contextu urządzeniach. 
Dalej w przykładowym kodzie stworzona jest kolajka \textbf{clCreateCommandQueue} kolejka powstaje w ramach contextu na konkretny device. Później na tym obiekcie kolejkowane będą zadania takie jak transfery pamięci czy wykonywane funkcje. W zależności czy kolejka zostanie stworzona z flagą CL\_QUEUE\_OUT\_OF\_ORDER\_EXEC\_MODE\_ENABLE lub bez niej, zadania te będą mogły być wykonywane równolegle, lub jedno po drugim w kolejności dodania do kolejki.
Kolejnym krokiem w powyższym kodzie jest stworzenie obiektu programu używając \textbf{clCreateProgramWithSource} lub \textbf{clCreateProgramWithBinary} pierwsza stworzy program zawierający nie skompilowany kod w języku OpenCL C, natomiast druga stworzy obiekt z binarnej wersji, wcześniej skompilowanej. Do stworzenia programu ze źródeł przekazywany jest ciąg znaków zawierający kernele, czyli funkcje które mogą zostać wykonane na urządzeniu.
Do wykonania inkrementacji każdego elementu bufora w przykładzie użyty zostanie następujący kernel
\lstinputlisting{listings/increment.cl}
W funkcji tej został pobrany unikalny numer aktualnie wykonywanego kernela w ramach globalnej work grupy, następnie element bufora pod tym indeksem jest inkrementowany. 
Po stworzeniu programu zawierającego kernele, należy zawołać \textbf{clBuildProgram} by kod kerneli w języku OpenCL C został skompilowany dla wskazanego urządzenia. W przypadku stworzenia programu ze źródeł w formie binarnej, tego kroku się nie wykonuje.
Następnym wykonanym krokiem jest stworzeniem obiektu bufora poprzez \textbf{clCreateBuffer}. Stworzony został obiekt reprezentujący obszar pamięci o podanym rozmiarze, który może być wykorzystany przy wykonywaniu kernela. W podanym kodzie powstanie bufor o rozmiarze 4096 bajtow czyli 1024 elementów typu int.
Następnie zostaje wykonane \textbf{clEnqueueMapBuffer} funkcja ta mapuje konkretny bufor na obszar pamięci dostępny z poziomu aplikacji. W tym przypadku w pamięci pod zwróconym wskaźnikiem ustawiamy w każdym bajcie wartość 13.
Aby przesłać pamięć z powrotem do obiektu bufora dostępnego z poziomu urządzenia na którym będzie wykonywany kernel wołamy \textbf{clEnqueueUnmapMemObject}.
Tak przygotowany bufor z danymi możemy ustawić jako argument kernela wołając \textbf{clSetKernelArg} podając w argumentach kernel, który w którym chcemy ustawić argument, index argumentu, jego typ oraz wskaźnik na obiekt który będzie argumentem funkcji.
Funkcja która uruchomi wykonanie kernela na wskazanym wcześniej urządzeniu jest \textbf{clEnqueueNDRangeKernel}, która umieści w kolejce do wykonywania wskazany kernel. W tym momencie podane jest także w ilu wymiarach odbędzie sie wykonywanie, podana jest wielkość lokalnej i globalnej work grupy.
Na samym końcu zawłoane zostaje \textbf{clFinish}, jest to funkcja blokująca po wykonaniu której mamy pewność, że wszystkie zakolejkowane na konkretnej kolejce operacje zostały wykonane.
W wyniku działania takiego kodu w buforze znajdować się będzie 1024 elementy o wartości 13131314.
\subsection{Możliwości i ograniczenia sprzętowe}
Każde urządzenie posiada ograniczenia związane ze specyfiką implementacji sterownika oraz barkiem zasobów sprzętowych. By dowiedzieć się jakie są maksymalne dostępne wartosci np dotyczące rozmiaru pamięci, ilości poszczególnych obiektów w kernelu czy maksymalnej liczbie dostępnych work itemów w ramach lokalnej work grupy możemy odpytać sterownik wołając clGetDeviceInfo podając konkretny parametr. Dzięki temu wykonywana aplikacja może dostosować się do ograniczeń sprzętowych. Oto wynik działania aplikacji "clinfo", która wypisuje wszystkie dostępne informacje o urządzeniu. W tym wypadku jest to telefon Xiaomi Mi a2 lite z procesorem graficznym Adreno 506
 \lstinputlisting{listings/clinfo.log}
 \subsection{Pomiary czasu wykonywania kerneli}
 Standard OpenCL zapewnia mechanizm do odczytywania stempli czasu z poszczególnych etapów wykonywania kernela. Służy do tego obiekt typu cl\_event stworzony przez funkcje clCreateUserEvent. Obiekt taki może zostać przekazany jako argument funkcji, która zostanie wykonana na urządzeniu np clEnqueueNDRangeKernel czy clWnqueueWriteBuffer.
 Po wykonaniu kernela z obiektu eventa można odczytać stemple czasu z jego wykonania. By wydobyć wartości należy zawołać clGetEventProfilingInfo przekazując jako argument obiekt eventu oraz jeden z czterech parametrów 
 \begin{itemize}
     \item CL\_PROFILING\_COMMAND\_QUEUED wartość opisuje czas urządzenia w którym komenda została dodana do kolejki.
     \item CL\_PROFILING\_COMMAND\_SUBMIT wartośc opisuje czas urządzenia w którym komenda została wysłana do urządzenia na którym zostanie wykonana.
     \item CL\_PROFILING\_COMMAND\_START wartość opisuje czas urządzenia w którym rozpoczęte zostało wykonywanie komenty na urządzeniu.
     \item CL\_PROFILING\_COMMAND\_END wartość opisuje czas urządzenia w którym wykonywanie komendy zostaje zakończone.
\end{itemize}
\subsection [Współdzielenie obiektów pomiędzy OpenCL i OpenGL]{Współdzielenie obiektów pomiędzy OpenCL i OpenGL}
OpenGL czyli Open Graphics Library to standard, ktory tak jak OpenCL jest zdefiniowany i utrzymywany przez grupę Khronos. Jest on wykorzystywany głównie do renderowania obiektów graficznych.
Istnieje możliwość stworzenia textury w OpenGL i użycie tego w kernelu OpenCL. Nie wszystkie urządzenia wspierają współdzielenie zasobów, urządzenie wspiera wspoldzielenie jeśli w liście rozszerzeń znajduje się cl\_khr\_gl\_sharing. By móc skorzystać z tej możliwości, context w OpenCL musi zostać stworzony w oparciu o wcześniej stworzony context OpenGL. Przy tworzeniu contextu podajemy strukturę, w której umieszczone są wskaźniki na kontekst w OpenGL i uchwyt na EGLDisplay. 
 \lstinputlisting{listings/contextProperties.txt}.
 Posiadając tak stworzony kontekst, możemy stworzyć obiekt obrazka w OpenCL w oparciu o stworzoną teksture w OpenGL
 \lstinputlisting{listings/sharedImage}. gdzie textureId to identyfikator tekstury, ktora chcemy współdzielić z openGL. 
 Przed wykorzystaniem współdzielonego obrazka należy zawołać clEnqueueAcquireGLObjects wskazując kolejkę, która będzie używać obrazka. Jest to wykonywane w celu wywłaszczenia obiektu przez OpenCL. Po zakończeniu operacji na obiekcie należy go zwolnić poleceniem clEnqueueReleaseGLObjects.
 \newpage
\section[Specyfika testowanych Urządzeń]{Specyfika testowanych Urządzeń}

 \newpage
\section[Testy Performancowe OpenCL na Mobilkach]{Testy Performancowe OpenCL na Mobilkach}

\subsection[Pierwsza sekcja]{Pierwsza sekcja}
\subsubsection{Podsekcja}
\lipsum[6]

\begin{figure}[H]
	\includegraphics[scale=0.8]{imgs/eti.png}
	\caption{Logo WETI}
\end{figure}

\lipsum[3]

\begin{figure}[H]
	\includegraphics[scale=0.8]{imgs/pg.jpg}
	\caption{Logo PG}
\end{figure}

\lipsum[7]

\subsubsection{I kolejna podsekcja}
\lipsum[2]

\begin{table}
\caption{Przykład krótkiej tabeli}
\label{tabela_1}
\begin{tabu} to \textwidth {| p{2cm} | p{2cm} | X |}
\hline
\rowcolor{lightgray} Nagłówek 1 & Nagłówek 2 & Nagłówek 3 \\ \hline
1 & lubię plaki & \lipsum[2] \\ \hline
2 & a ja nie & \lipsum[2] \\ \hline
\end{tabu}
\end{table}

\lipsum[3]

\subsection[Druga sekcja]{Druga sekcja}
\subsubsection{Podsekcja sekcji drugiej}
\lipsum[4]

\begin{longtabu} to \textwidth {| p{2cm} | X |}
\caption{Przykład długiej tabeli} \label{tabela_2} \\[-2mm] %-2mm jest wymagane, poniewaz caption w longtabu jest uznawany jako wiersz i działa na niego ustawione ograniczenie min. 2mm odstępu przed następnym wierszem
\hline
\rowcolor{lightgray} Nagłówek 1 & Nagłówek 2 \\ \hline
\endfirsthead

\rowcolor{lightgray} Nagłówek 1 & Nagłówek 2 \\ \hline
\endhead

1 & \lipsum[5] \\ \hline
2 & \lipsum[5] \\ \hline
3 & \lipsum[5] \\ \hline
4 & \lipsum[5] \\ \hline
5 & \lipsum[5] \\ \hline
6 & \lipsum[5] \\ \hline
\end{longtabu}

\lipsum[3]

\subsubsection{Kod}
\lipsum[2]

\lstinputlisting{listings/main.cpp}

\lipsum[2]
 \newpage
\section[Aplikacje wykorzystujące OpenCl Przypadki Użycia]{Aplikacje wykorzystujące OpenCl Przypadki Użycia}

\subsection[Pierwsza sekcja]{Pierwsza sekcja}
\subsubsection{Podsekcja}
\lipsum[6]

\begin{figure}[H]
	\includegraphics[scale=0.8]{imgs/eti.png}
	\caption{Logo WETI}
\end{figure}

\lipsum[3]

\begin{figure}[H]
	\includegraphics[scale=0.8]{imgs/pg.jpg}
	\caption{Logo PG}
\end{figure}

\lipsum[7]

\subsubsection{I kolejna podsekcja}
\lipsum[2]

\begin{table}
\caption{Przykład krótkiej tabeli}
\label{tabela_1}
\begin{tabu} to \textwidth {| p{2cm} | p{2cm} | X |}
\hline
\rowcolor{lightgray} Nagłówek 1 & Nagłówek 2 & Nagłówek 3 \\ \hline
1 & lubię plaki & \lipsum[2] \\ \hline
2 & a ja nie & \lipsum[2] \\ \hline
\end{tabu}
\end{table}

\lipsum[3]

\subsection[Druga sekcja]{Druga sekcja}
\subsubsection{Podsekcja sekcji drugiej}
\lipsum[4]

\begin{longtabu} to \textwidth {| p{2cm} | X |}
\caption{Przykład długiej tabeli} \label{tabela_2} \\[-2mm] %-2mm jest wymagane, poniewaz caption w longtabu jest uznawany jako wiersz i działa na niego ustawione ograniczenie min. 2mm odstępu przed następnym wierszem
\hline
\rowcolor{lightgray} Nagłówek 1 & Nagłówek 2 \\ \hline
\endfirsthead

\rowcolor{lightgray} Nagłówek 1 & Nagłówek 2 \\ \hline
\endhead

1 & \lipsum[5] \\ \hline
2 & \lipsum[5] \\ \hline
3 & \lipsum[5] \\ \hline
4 & \lipsum[5] \\ \hline
5 & \lipsum[5] \\ \hline
6 & \lipsum[5] \\ \hline
\end{longtabu}

\lipsum[3]

\subsubsection{Kod}
\lipsum[2]

\lstinputlisting{listings/main.cpp}

\lipsum[2]
 \newpage
\section[Aplikacje wykorzystujące OpenCl Przypadki Użycia]{Aplikacje wykorzystujące OpenCl Przypadki Użycia}
OpenCL jest biblioteką, która z pewnością może znaleźć zastosowanie w wielu aplikacjach z systemem Android. Brak oficjalnego wsparcia api na Androidzie sprawia, że ilość informacji o używaniu tej biblioteki przez programy na platforme Android jest znikoma. O używaniu tego api przez aplikacje można znaleźć tylkokilka wzmianek w internecie. W poniższych podrozdzialach opisano wykożystanie tej biblioteki w kilku aplikacjach.
\subsection[Tensorflow]{Tensorflow}
Tensorflow to biblioteka, wykorzystywana do uczenia maszynowego. Tensorflow jest często wykorzystywany do treningów głębokich sieci neuronowych. Proces uczenia sieci neuronowych, wymaga wielokrotnego przemnażania wielu macierzy. Samo wykorzystanie wyuczonego modelu również wiąże się z procesem mnożenia macierzy. Tensorflow to oprogramowanie potrafiące wykożystać różne api służące do obliczeni, takie jak Cuda, Sycl czy OpenCL. W przypadku wersji programu na system Android głównym api używanym przez tensorflow jest OpenCL. Jak wynika z artykułu \cite{Tensor}, użycie api OpenCL umożliwia poprawienie wydajności o 100\% w stosunku do uzywania OpenGL. Przy optymalizacji programu używająć OpenCL pomocne były narzędzia dostarczone przez producenta GPU Adreno. Poniewaz api OpenCL nie jest oficjalnie wspierane, aplikacja tensorflow na starcie odpytuje urządzenie o dostępność biblioteki. W przypadku gdy biblioteka ta jest nie dostępna ładowany jest silnik działający w oparciu o api OpenGL.
\subsection[Adobe]{Adobe}
\subsection[Testy mierzące wydajność i jakość urządzenia]{Testy mierzące wydajność i jakość urządzenia}
 \newpage
\section[Podsumowanie]{Podsumowanie}
Celem niniejszej pracy była analiza środowiska jakim jest OpenCL na platformach mobilnych. Sprawdzono, że z pośród najpopularniejszych systemów mobilnych, tylko Android posiada możliwość obsługi tego api. Przy realizacji zadania sprawdzono, w jaki sposób programy, mogą łączyć się z biblioteką implementującą interfejs i jak to wpływa na ich przenośność. W ramach pracy skonstruowano i wykonano kilka testów mających za zadanie przetestować możliwości urządzeń z systemem Android w pracy z biblioteką OpenCL, pod kątem wydajnościowym, jak i współpracy api OpenCL z OpenGL przy wyświetlaniu obrazu z kamery urządzenia. Testy zostały przeprowadzone na urządzeniach opisanych i porównanych w niniejszym dokumencie. Wykonane aplikacje zostały dokładnie opisane, a ich wyniki w formie wykresów zostały przedstawione w pracy i przeanalizowane. W ramach pracy został wykonane poszukiwanie aplikacji wykorzystujących api OpenCL.

Przy realizacji pracy, problemem okazał się brak oficjalnego wsparcia dla OpenCL przez system Android. Pomimo to producenci procesorów graficznych na urządzenia mobilne, tworzą oprogramowanie, które implementuje funkcjonalość api OpenCL. By móc skorzystać z biblioteki należy albo bezpośrednio połączyć aplikację z konkretną dynamiczną biblioteką, albo ładować ją w trakcie wykonywania kodu. Przy zastosowaniu pierwszego rozwiązania aplikacja może być wykonywana tylko na urządzeniu z którego pochodzi biblioteka, na stałe połączona z programem. W drugim przypadku biblioteka jest ładowana dynamicznie. Interfejs OpenCL na platformie Android dostępny jest jedynie z poziomu natywnego kodu w C++, by używać jej w kodzie aplikacji, niezbędne jest wykonywanie części OpenCL po stronie natywnej biblioteki, lub wykonanie dodatkowej biblioteki, która będzie nakładką na funkcje w C++. 

W ramach przeprowadzonych testów sprawdzono jaką moc obliczeniową mają poszczególne urządzenia przy użyciu api OpenCL. Wyraźnie zaobserwowano, że urządzenia wyprodukowane przez Adreno, reagują podobnie na zmianę typu danych i w raz z używaniem większych typów wektorowych ich moc obliczeniowa rosła, natomiast w przypadku procesorów Mali przy zastosowaniu typów wektorowych moc obliczeniowa pozostawała na tym samym poziomie. Procesory Mali są w stanie wykonywać więcej obliczeń na typach całkowitych, natomiast układy Adreno lepiej radza sobie z operacjami na typach zmiennoprzecinkowych. Zweryfikowano, że rozmiary lokalnych work grup znacząco wpływają na czas wykonywania mnożenia macierzy, które wymaga czytania dalekich od siebie komórek pamięci. Urządzenia najlepsze czasy osiągały gdy rozmiary grup w wymiarze X i Y są równomierne. W przypadku takich grup podczas wykonywania wątki w ramach simd, mogły dostępować elementy z pamięci podręcznej. W takim przypadku kosztowne transfery alokacji do pamięci cache występują rzadziej. W tym teście kluczową rolą jest szybki dostęp do pamięci, dlatego im urządzenia miały lepszy typ pamięci, tym lepiej radziły sobie z tym zadaniem. 

W ramach pracy, sprawdzono możliwości urządzeń do modyfikacji podglądu z kamery przy pomocy OpenCL. Okazało się, że podgląd z kamery dostarczany jest w formacie NV21, obsługi którego api OpenCL nie definiuje. Dlatego by móc użyć obrazu z kamery w wykonywanym kernelu, należało przetransformować go do formatu RGBA. Konwersja znacząco wpłynęła jedynie na urządzenie z procesorem graficznym Mali. Na urządzeniach z GPU od Adreno proces zamiany obrazka na format RGBA i wpisanie go do wyświetlanej textury nie wpłynął na ilość wyświetlanych klatek na sekundę. Po nałożeniu prostych filtrów jak zamianę na obraz w skali szarości lub filtr max rgb, liczba prezentowanych obrazów na sekundę nie zmieniła się. Natomiast użycie filtru uśredniającego, w którym do wyliczenia wartości jednego piksela potrzebne było odczytanie i uśrednienie wartości 25 pikseli, znacząco obniża ilość wyświetlanych klatek, dla niektórych urządzeń nawet o ponad połowę.

Istnieje mało informacji na temat korzystania z api OpenCL w aplikacjach użytkowych. Każdy z producentów procesorów graficznych dla urządzeń z systemem Android, posiada implementację takiego sterownika, co pozwala przypuszczać, że jest on częściej używany niż tylko do testów sprawdzających możliwości urządzenia. Problem z dotarciem do informacji na temat używania tego api przez aplikacje najprawdopodobniej wynika z tego, że twórcy programów zazwyczaj nie umieszczają nigdzie listy wszystkich użytych api, frameworków, czy wersji języka programowania.

W przypadku kontynuowaniu pracy nad tematem, można by było znaleźć sposób na sprawdzenie zużycia energii przy wykonywaniu kerneli. Innym kierunkiem kontynuowania prac mogłoby być sprawdzenie jak OpenCL na platformie Android działa na urządzeniach z kartami graficznymi Intela Amd czy Nvidi.
 \newpage

\renewcommand{\baselinestretch}{1.0}\normalsize	% żeby w wykazach była pojedyncza interlinia

\addcontentsline{toc}{section}{Wykaz literatury}
\nocite{*}
\printbibliography
\newpage

\addcontentsline{toc}{section}{\listfigurename}
\listoffigures
\newpage

\addcontentsline{toc}{section}{\listtablename}
\listoftables

\renewcommand{\baselinestretch}{1.5}\normalsize	% powrót do interlinii 1.5 na wypadek dodatków

% nie dodawaj obrazków i tabel z dodatków do list powyżej
\let\svaddcontentsline\addcontentsline
\renewcommand\addcontentsline[3]{%
  \ifthenelse{\equal{#1}{lof}}{}%
  {\ifthenelse{\equal{#1}{lot}}{}{\svaddcontentsline{#1}{#2}{#3}}}}

\newpage

\begin{appendices}
\addtocontents{toc}{\protect\setcounter{tocdepth}{1}}

\section{Przykładowy dodatek} \label{app:example}

\subsection{Sekcja}
 \newpage
\end{appendices}

\end{document}	% musi być na samiutkim końcu
